\section{Начини за споделяне на един порт}
\label{sec:sharing}

Задачата „един порт, няколко протокола“ не е нова. Подходите се различават по това къде
се взема решението, колко байта са нужни за него и кой процес провежда TLS
ръкостискането. Този раздел разглежда четири подхода. За двата инструмента от трети
страни се цитира само това, което пише в документацията им.

\subsection{Отделен порт за всеки протокол}

Това е традиционният подход: 80 за некриптиран HTTP и 443 за HTTPS \cite{rfc9110}.
Решение не се взема, защото портът вече е решението. Препоръките на IETF подкрепят
отделния порт само за TLS пред динамичното надграждане на некриптирана връзка до TLS
\cite[раздел 3.2]{rfc9325}. Причината е, че при отделния порт крайното състояние на
канала е известно още преди първия байт.

\subsection{Надграждане на некриптирана връзка}

Някои протоколи започват некриптирано и по-късно преминават към TLS с команда, например
STARTTLS при пощенските протоколи. RFC 9325 изброява този подход и предупреждава, че
началото на връзката не е защитено и нападател може да го промени така, че
надграждането да не се случи \cite[раздел 3.2]{rfc9325}. Атаката срещу командата
STARTTLS е описана и в RFC 7457 \cite[раздел 2.2]{rfc7457}.

\subsection{Опцията port-share на OpenVPN}

OpenVPN може да работи като TCP сървър. С опцията \code{-{}-port-share} той споделя порта
си с друго приложение, например HTTPS сървър \cite{openvpn26}. Когато OpenVPN разпознае,
че връзката не е по неговия протокол, той я препраща като посредник към зададените адрес
и порт. По желание той записва във временни файлове адреса и порта на източника и на
целта за всяка препратена връзка. Документацията изрично посочва, че опцията не е
реализирана за Windows.

Тук решението се взема в процеса на OpenVPN. TLS ръкостискането на уеб трафика се
провежда от уеб сървъра зад него, а не от OpenVPN. Цената е допълнителен скок: всяка
некриптирана или HTTPS връзка минава през посредник, който копира байтовете между два
сокета.

\subsection{sslh}

sslh е самостоятелен мултиплексор. Той приема връзки на зададени портове и ги препраща
нататък според проверки върху първия пакет данни, изпратен от клиента \cite{sslh}.
Документацията му изброява проверки за HTTP, TLS (включително по SNI и ALPN), SSH,
WireGuard, OpenVPN, tinc, XMPP и SOCKS5, както и за всеки протокол, който може да се
опише с регулярен израз. Както при OpenVPN, TLS се провежда от услугата зад
мултиплексора, а не от него самия.

\subsection{Сравнение}

\begin{table}[H]
\centering
\caption{Сравнение на подходите за споделяне на порт}
\label{tab:sharing}
\begin{tabularx}{\textwidth}{@{}L{2.9cm}YYY@{}}
\toprule
Подход & Къде се решава & Кой провежда TLS & Допълнителен скок \\
\midrule
Отделни портове & никъде, портът е решението & сървърът на порт 443 & няма \\
OpenVPN port-share & в процеса на OpenVPN & уеб сървърът зад него & да, посредник \\
sslh & в отделен процес, по първия пакет & услугата зад него & да, препращане \\
\code{octet} & в самия сървър, по първия октет & самият сървър & няма \\
\bottomrule
\end{tabularx}
\end{table}

Колко точно октета четат OpenVPN и sslh, за да вземат решение, не е посочено в цитираната
документация и работата не прави предположения за това. Сравнението показва две неща.
Първо, \code{octet} е единственият от разгледаните подходи, при който процесът, който
разпознава протокола, е и процесът, който провежда ръкостискането. Няма посредник и
няма втори сокет. Второ, OpenVPN и sslh различават много протоколи, а \code{octet} само
два. Именно затова на \code{octet} му стига един октет. Мултиплексор, който трябва да
различи SSH от OpenVPN от HTTP, не може да си позволи това.
